\section{Introduction}\label{sec:introduction-compsuite}

% add some context of AI and how it can be used to solve the problem of library incompatibility issues.
% add some context of benchmarks

Modern software systems are becoming increasingly complex due to the
need for integrating various components developed by different teams
or organizations. These components are often subject to continuous
evolution, and as a result, ensuring that new
upgrades to third-party libraries do not cause any compatibility
issues with the existing software system is a challenging task. The
complexity of these systems and the number of dependencies involved
make it difficult to anticipate and identify incompatibilities that
may arise from updates to external components. Incompatibility issues resulting from upgrades to external components can
compromise the reliability of software systems, potentially leading to
significant financial losses for the organizations that rely on these systems.

Many techniques have been proposed to address third-party library
compatibility issues, including regression
testing~\cite{mezzetti2018type,moller2019model}, static
analysis~\cite{foo2018efficient}, dependency conflict
detection~\cite{wang2021will}, and client-specific compatibility
checking~\cite{mora2018client,zhu2019framework}. These techniques address
library compatibility issues in different dimensions and have been
evaluated with their own isolated datasets.

In recent years, Large Language Models (LLMs) have achieved remarkable success across numerous automated software engineering tasks. At the heart of this progress lie high-quality open-source datasets, which serve as the cornerstone of model capabilities. These datasets provide essential domain-specific training material for learning real-world complex software patterns. Furthermore, the precision of their annotations and the reproducibility of their instances directly impact the performance upper bound of AI tools on various software engineering tasks. For example, \textit{Defects4J}~\cite{just2014defects4j} contains over 800 reproducible bugs drawn from 17 open-source Java projects, enabling evaluation of whether AI can generate patches that pass test suites; \textit{CoRe}~\cite{xie2025core} includes more than 10,000 task instances covering data dependency, control dependency, and information flow across multiple programming languages for evaluating the ability of fundamental static analysis; \textit{SWE-bench}~\cite{jimenez2023swe} provides an end-to-end benchmark comprising a large collection of real GitHub issues, designed to assess the end-to-end capability of AI in resolving software engineering problems.

However, these datasets mainly focus on internal problems of projects, leaving largely uncovered risks introduced by external changes. Although existing LLM-based tools have demonstrated impressive performance in code comprehension, they face substantial challenges in understanding and reasoning about compatibility issues induced by library upgrades, severely limiting their capacity to manage external evolutionary risks. This limitation stems from two interrelated factors. First, library upgrade compatibility issues are inherently complex, involving multiple components and intricate dependency relationships. Current LLM reasoning is often insufficient to model cross-component interactions or infer the impact of changes. Second, such issues are difficult to detect and some of them may only trigger under specific environments and test cases. Additionally, constructing reproducible problem instances requires non-trivial configuration, which further exacerbates the scarcity of data. Consequently, a structured, high-quality, and fully reproducible dataset of real-world library upgrade incompatibilities is essential for enhancing and evaluating LLM-based tools in this domain.


An excellent dataset has the potential to serve as a valuable reference for future research in
this field. However, composing the dataset requires intricate
manual validation, e.g., confirming whether the cause of a test
failure is due to runtime exception, assertion violation, or other reasons.
Therefore, we propose \toolcompsuite, the first incompatibility issue dataset focusing on
library behavioral incompatibility with concrete reproducible test cases.
Each test case is isolated and validated, enabling the direct manifestation of the
incompatibilities.


An excellent dataset has the potential to serve as a valuable reference for future research in
this field. However, composing the dataset requires intricate
manual validation, e.g., confirming whether the cause of a test
failure is due to runtime exception, assertion violation, or other reasons.
Therefore, we propose \toolcompsuite, the first incompatibility issue dataset focusing on
library behavioral incompatibility with concrete reproducible test cases.
Each test case is isolated and validated, enabling the direct manifestation of the
incompatibilities.


\toolcompsuite comprises 123 real-world Java client-library pairs such that upgrading any library results in
incompatibility issues for the corresponding client.
Every incompatibility issue in \toolcompsuite contains a test case created by developers, allowing for the
reproduction of the issue.
On top of this dataset, we also developed an automated command-line interface, which streamlines
all processes of the reproduction, such as downloading and compiling a projects, running target
tests and re-runing the tests after a library upgrade.
With this infrastructure, users may reproduce an incompatibility issue programmatically with
minimal efforts.




% In this study, we present \toolcompsuite, a dataset comprising 123 real-world Java client-library pairs in which upgrading the library results in compatibility problems for the respective client. Each compatibility issue in CompSuite is linked to a test case created by developers, allowing for the reproduction of the issue. Additionally, the dataset features an automated tool that streamlines the execution and verification of each problem. This setup enables users to swiftly assess any compatibility issue with a single click or to reproduce the issue methodically for an in-depth examination.


\textbf{Contribution}
To summarize, we make the following contributions in this chapter:%
\begin{enumerate}[leftmargin=*,topsep=2pt]
\item We construct a dataset, \toolcompsuite, including 123
  reproducible, real-world client-library pairs that manifest
  incompatibility issues when upgrading the library. These data points
  originate from 88 clients and 104 libraries.
  %% We forked each project and identified the base version that
  %% passed all tests before encountering incompatibility
  %% issues. Based on that version, we performed library updates and
  %% utilized a version control system to create isolated commits,
  %% with each version differing only in the library version.
\item We created an automated command-line interface for the dataset.
With this interface, users are able to programmatically replicate an incompatibility issue
from the dataset with a single command.
The interface also offers separate commands for each step involved in the reproduction of
incompatibility issues.
\end{enumerate}


% \todo{TODO: Add a paragraph here: We envision that \toolcompsuite can power various
%   program analysis techniques including xxx, xxx, and xxx. Briefly
%   describe how each of these techniques would use \toolcompsuite.}

We envision that \toolcompsuite to be used to evaluate various program analysis techniques, including
compatibility checking, module-level regression testing selection, and debugging techniques.
%For example, compatibility checkers and detecters can use \toolcompsuite as a performance evaluator, while
%module-level RTS tools utilize it as a safety evaluator to verify whether all relevant test cases
%have been selected. Additionally, \toolcompsuite is able to assist debugging techniques to do better root
%cause analysis.
Furthermore, this is also very high-quality and valuable data for AI, which can help models better understand such type of issue.
More detailed information can be found in \sref{sec:app-compsuite}.

The dataset and tool are available at: \url{https://github.com/compsuite-team/compsuite}.

